\section{The MPME sequence and its forward model}\label{sec:model}

\subsection{Sequence}

The MPME sequence~\cite{cheng2019} is a 3D gradient-echo sequence without RF spoiling: the RF
phase is constant and the repetition time $\TR$ is fixed. Only the readout gradient $G_x$ is
unbalanced; its net area per repetition, $A=\int_0^{\TR}G_x\,dt$, dephases the magnetisation
by many cycles across a voxel (the \emph{fully dephased} regime). Configuration states of
different order are therefore displaced in $k_x$ by multiples of $\gamma A$, and a state of
order $k$ created by an RF pulse is refocused when the moment accumulated since the pulse
equals $-kA$.

In the published design each of three readout windows of alternating polarity sweeps the
moment across three adjacent multiples of $A$ (lobe areas $+1.5:-3:+3:-3:+0.5$ in units of
$-A$), so that every window samples the pathways $k\in\{+1,0,-1\}$, separated by the time
needed to traverse one $A$, which is $1/\mathrm{BW}\approx2$\,ms (\cref{fig:sequence}). The
visiting order reverses in alternate windows. Two acquisitions with the same gradients but
very different flip angles are performed; the second covers only the central part of
$k$-space (\cref{tab:protocol}). The actual flip angles are
\begin{equation}
  \alpha_i=\Bone\,\alpha_i^{\rm nom},
  \label{eq:b1}
\end{equation}
so $\Bone$ scales all pulses by the same factor.

\begin{figure}[!htbp]
\centering
\includegraphics[width=\linewidth]{sequence_timing.png}
\caption{One repetition of the $[+1,0,-1]$ scheme. Top: readout gradient (prephaser, three
readout windows of alternating polarity, rewinder); net area $A$ per repetition. Bottom: the
moment accumulated since the RF pulse; pathway $k$ is refocused (circles) whenever $m(t)=-k$.
The time from the pulse to the first echo (4.5\,ms) is an assumption; the paper does not
state it.}
\label{fig:sequence}
\end{figure}

\begin{table}[!htbp]
\centering
\caption{Published in-vivo protocol~\cite{cheng2019} and the values used in this note.}
\label{tab:protocol}
\small
\begin{tabular}{ll}
\toprule
Repetition time & $\TR=25$\,ms (both scans)\\
Nominal flip angles & $\alpha_1^{\rm nom}=15^\circ$, $\alpha_2^{\rm nom}=330^\circ$ (hard pulses)\\
Pathways & $k\in\{+1,0,-1\}$; three readout windows\\
Pathway spacing & $2.0$\,ms in vivo ($1.86$\,ms phantom); 501\,Hz/pixel\\
Echo times (ms) & $k=+1$: $4.5,14.5,16.5$;\ $k=0$: $6.5,12.5,18.5$;\ $k=-1$: $8.5,10.5,20.5$\\
Scan 2 coverage & central $25\%\times25\%$ of $k_y$--$k_z$\\
Resolution, matrix & $1.2\times1.0\times1.2$\,mm$^3$, $160\times192\times160$, 11.5\,min\\
\bottomrule
\end{tabular}
\end{table}

\subsection{Configuration states}

Let $\psi\in[0,2\pi)$ be the intravoxel phase imparted by one repetition of the unbalanced
gradient and expand
\begin{equation}
  M^+(\psi)=\sum_kF_k\,\E^{ik\psi},\qquad M_z(\psi)=\sum_kZ_k\,\E^{ik\psi},\qquad
  Z_{-k}=\overline{Z_k}.
\end{equation}
One repetition consists of~\cite{hennig1991,weigel2015}:
\begin{align}
  \text{RF:}\quad F_k'&=\cos^2\tfrac{\alpha}{2}F_k+\E^{2i\varphi}\sin^2\tfrac{\alpha}{2}\overline{F_{-k}}
     -i\E^{i\varphi}\sin\alpha\,Z_k,\nonumber\\
  Z_k'&=-\tfrac{i}{2}\E^{-i\varphi}\sin\alpha\,F_k+\tfrac{i}{2}\E^{i\varphi}\sin\alpha\,\overline{F_{-k}}
     +\cos\alpha\,Z_k,\\
  \text{relaxation:}\quad F_k&\to E_2F_k,\quad Z_k\to E_1Z_k\ (k\ne0),\quad
     Z_0\to E_1Z_0+\Mzero(1-E_1),\nonumber\\
  \text{gradient:}\quad F_k&\to F_{k+1}.\nonumber
\end{align}
We write $a_k(\alpha,\Tone,\Ttwo,\dw)$ for the steady-state amplitude of $F_k$ immediately
after the pulse, per unit $\Mzero$.

\subsection{Exact steady state}

The steady state can be computed without truncating the configuration space. An isochromat
whose total precession per repetition is $\beta=\psi+\dw\TR$ has the post-pulse steady state
\begin{equation}
  \bigl(I-R\,R_z(\beta)\,\mathrm{diag}(E_2,E_2,E_1)\bigr)\,\mathbf M^+=(1-E_1)\,R\,\hat{\mathbf z},
  \qquad R=R_z(\varphi)R_x(\alpha)R_z(-\varphi),
  \label{eq:iso}
\end{equation}
and $a_k$ is the $k$-th Fourier coefficient of $M^+(\psi)=M_x^++iM_y^+$ over $\psi$. On a uniform
grid of $n$ isochromats this is a discrete Fourier transform; orders with $|k|<n/2$ are exact
up to aliasing of orders differing by $n$, which decay as $E_2^{n}$-like factors. For every
tissue in this note, including CSF ($\Ttwo=1.5$\,s) at $330^\circ$, $n=128$ reproduces
$n=16\,384$ to $10^{-15}$.

\begin{proposition}[Off-resonance only adds phase]\label{prop:offres}
In the fully dephased regime
\[
  a_k(\dw)=a_k(0)\,\E^{ik\dw\TR};
\]
in particular $|a_k|$ does not depend on $\dw$.
\end{proposition}
\begin{proof}
By \eqref{eq:iso} the steady state depends on $\psi$ and $\dw$ only through $\beta=\psi+\dw\TR$,
so $M^+_{\dw}(\psi)=M^+_0(\psi+\dw\TR)$. Shifting a $2\pi$-periodic function by $\theta$ multiplies
its $k$-th Fourier coefficient by $\E^{ik\theta}$.
\end{proof}

\subsection{Two invariants per scan}

\begin{theorem}[Exact two-invariant reduction]\label{thm:uv}
Let $\varphi=0$ and, for one scan with flip angle $\alpha$, define
\begin{equation}
  u=\frac{E_1-\cos\alpha}{1-E_1\cos\alpha},\qquad
  v=\frac{\sin\alpha\,(1-E_1)}{1-E_1\cos\alpha}.
  \label{eq:uv}
\end{equation}
Then for every pathway $k$ and every $\dw$
\begin{equation}
  a_k(\alpha,\Tone,\Ttwo,\dw)=v\;G_k(u,E_2,\dw),
  \label{eq:uvform}
\end{equation}
where $G_k$ does not depend on $\alpha$ or $\Tone$ otherwise.
\end{theorem}
\begin{proof}
Let $(x,y,z)$ be the pre-pulse state of one isochromat and $w=x+iy$. The pulse $R_x(\alpha)$
maps it to $(x,\ cy-sz,\ sy+cz)$ with $c=\cos\alpha$, $s=\sin\alpha$; relaxation and precession by
$\beta$ then give the next pre-pulse state. The steady state therefore satisfies
\begin{equation}
  w=E_2\E^{i\beta}\bigl(x+i(cy-sz)\bigr),\qquad z=E_1(sy+cz)+1-E_1 .
  \label{eq:steady2}
\end{equation}
The second equation gives $z=(E_1sy+1-E_1)/(1-E_1c)$; note $1-E_1c>0$. Substituting,
\[
  cy-sz=y\,\frac{c-E_1c^2-E_1s^2}{1-E_1c}-\frac{s(1-E_1)}{1-E_1c}=-u\,y-v .
\]
The post-pulse transverse magnetisation is $m_\perp=x+i(cy-sz)=x-iuy-iv$, and $w=E_2\E^{i\beta}m_\perp$.
Eliminating $x=\Re w$ and $y=\Im w$,
\begin{equation}
  m_\perp=\Re\bigl(E_2\E^{i\beta}m_\perp\bigr)-iu\,\Im\bigl(E_2\E^{i\beta}m_\perp\bigr)-iv .
  \label{eq:mperp_lin}
\end{equation}
This is a real-linear equation for $m_\perp$ (two real unknowns) whose coefficients involve only
$u$, $E_2$ and $\beta$, and whose inhomogeneous term is $-iv$. Its $2\times2$ real matrix is
$I-E_2\,D_u\,R(\beta)$ with $D_u=\mathrm{diag}(1,u)$ and $R(\beta)$ a planar rotation; since
$|u|\le1$ and $E_2<1$, $\|E_2D_uR(\beta)\|\le E_2<1$, so it is invertible. Hence
$m_\perp=v\,g(u,E_2,\beta)$ with $g$ independent of $\alpha$ and $\Tone$, and taking Fourier
coefficients over $\psi$ (with $\beta=\psi+\dw\TR$) gives \eqref{eq:uvform}.
\end{proof}

\begin{lemma}\label{lem:absv}
$|u|\le1$ and $|v|=\sqrt{1-u^2}\,\sqrt{\tanh\bigl(\TR/(2\Tone)\bigr)}$.
\end{lemma}
\begin{proof}
From \eqref{eq:uv}, $1-u=(1-E_1)(1+c)/(1-E_1c)$ and $1+u=(1+E_1)(1-c)/(1-E_1c)$, both
non-negative, so $|u|\le1$ and $1-u^2=(1-E_1^2)\,s^2/(1-E_1c)^2$. Hence
$v^2=s^2(1-E_1)^2/(1-E_1c)^2=(1-u^2)(1-E_1)/(1+E_1)$, and $(1-E_1)/(1+E_1)=\tanh(\TR/2\Tone)$ for
$E_1=\E^{-\TR/\Tone}$.
\end{proof}

The invariants take their simplest form in terms of the Ernst angle $\alpha_E=\arccos E_1$, the
flip angle that maximises the spoiled-gradient-echo signal of the tissue.

\begin{proposition}[Ernst-angle form of the invariants]\label{prop:ernst}
Let $\zeta=\tan(\alpha_E/2)$ and, for a flip angle $\alpha$ that is not an odd multiple of $\pi$,
\begin{equation}
  \xi=\frac{\tan(\alpha/2)}{\tan(\alpha_E/2)},\qquad \vartheta=\arctan\xi\in(-\tfrac\pi2,\tfrac\pi2).
  \label{eq:xi}
\end{equation}
Then $\zeta^2=(1-E_1)/(1+E_1)=\tanh(\TR/2\Tone)$ and
\begin{equation}
  u=\frac{\xi^2-1}{\xi^2+1}=-\cos2\vartheta,\qquad
  v=\zeta\,\frac{2\xi}{1+\xi^2}=\zeta\sin2\vartheta .
  \label{eq:ernstform}
\end{equation}
\end{proposition}
\begin{proof}
Write $q=\tan(\alpha/2)$, so $\cos\alpha=(1-q^2)/(1+q^2)$ and $\sin\alpha=2q/(1+q^2)$, and likewise
$E_1=\cos\alpha_E=(1-\zeta^2)/(1+\zeta^2)$, which gives $\zeta^2=(1-E_1)/(1+E_1)=\tanh(\TR/2\Tone)$.
Then $E_1-\cos\alpha=2(q^2-\zeta^2)/D$, $1-E_1\cos\alpha=2(q^2+\zeta^2)/D$ and
$\sin\alpha\,(1-E_1)=4q\zeta^2/D$ with $D=(1+q^2)(1+\zeta^2)$. Substituting in \eqref{eq:uv},
$u=(q^2-\zeta^2)/(q^2+\zeta^2)$ and $v=2q\zeta^2/(q^2+\zeta^2)$; dividing numerator and denominator
by $\zeta^2$ gives \eqref{eq:ernstform}, and $\xi=\tan\vartheta$ gives the trigonometric forms.
\end{proof}

\Cref{lem:absv} is the identity $u^2+(v/\zeta)^2=1$: the pair $(u,v/\zeta)$ lies on the unit
circle at angle $2\vartheta$. $\xi$ measures where the actual flip angle lies relative to the
Ernst angle: $\xi<1$ ($u<0$) below it, $\xi=1$ ($u=0$, $v=\zeta$, the largest value of $v$) at
it, and $\xi\to\infty$ ($u\to1$, $v\to0$) as $\alpha\to\pi$. For $\alpha\ll1$ and $\TR\ll\Tone$,
\begin{equation}
  \xi\approx\frac{\alpha}{\alpha_E}\approx\alpha\sqrt{\frac{\Tone}{2\TR}},\qquad
  \zeta\approx\sqrt{\frac{\TR}{2\Tone}} ,
  \label{eq:xismall}
\end{equation}
so $\xi^2$ is the exact counterpart of the similarity variable $c=\alpha^2\Tone/2\TR$ of
\cref{prop:scaling}. At $15^\circ$ and $\TR=25$\,ms the Ernst angles of WM, GM and CSF are $13.8^\circ$,
$11.0^\circ$ and $6.4^\circ$, so $\xi=1.09$, $1.37$ and $2.36$ ($u=0.08$, $0.30$ and $0.69$). The
small-angle forms are close but not exact: for WM, $\xi^2=1.1787$, $(\alpha/\alpha_E)^2=1.1767$ and
$\alpha^2\Tone/2\TR=1.1652$.

On a logarithmic half-angle axis the invariants take an even simpler form. With
$s=\ln|\tan(\alpha/2)|$ and $s_E=\ln\zeta=\ln\tan(\alpha_E/2)$, so that $|\xi|=\E^{s-s_E}$, \eqref{eq:ernstform} reads
\begin{equation}
  u=\tanh(s-s_E),\qquad |v|=\zeta\,\operatorname{sech}(s-s_E),\qquad \sgn v=\sgn\tan(\alpha/2).
  \label{eq:logtan}
\end{equation}
As a function of $s$, the spoiled-gradient-echo amplitude $|v|$ is a bump of fixed shape
($\operatorname{sech}$), centred at the tissue's Ernst angle $s_E$ and of height $\zeta$. $\Tone$ sets the position and
height of the bump, and the actual flip angle sets the point $s$ at which a scan samples it.
Every log-derivative with respect to $\Bone$ follows from
$\mathrm d\ln|\tan(x/2)|/\mathrm d\ln x=x/\sin x$.

\begin{corollary}[One scan cannot separate $\Mzero$, $\Bone$, $\Tone$]\label{cor:onescan}
For a single scan, with $\Ttwo$, $\Ttwos$, $\dw$, $\phi_0$ fixed, the signal depends on
$(\Mzero,\Bone,\Tone)$ only through the two functions $(u,\ \Mzero v)$. The Jacobian of the
signal with respect to $(\log\Mzero,\log\Bone,\log\Tone)$ has rank at most two, and the
Fisher information about these three parameters is singular at every flip angle.
\end{corollary}
\begin{proof}
By \cref{thm:uv} and the echo factors below (which do not involve $\Mzero$, $\alpha$ or $\Tone$
except through the prefactor $\Mzero$), $S=\Mzero v\,h(u)$ for a function $h$ independent of
$(\Mzero,\Bone,\Tone)$. By the chain rule, the $3$ columns of $\partial S/\partial(\log\Mzero,
\log\Bone,\log\Tone)$ are linear combinations of the $2$ vectors $\partial S/\partial(u,\Mzero v)$,
so they span at most a two-dimensional space. The Fisher information restricted to these
parameters is $J^{\mathsf H}J/\sigma^2$ with $J$ of rank $\le2$, hence singular.
\end{proof}

\subsection{What the invariants mean, and what one scan measures}\label{sec:invariants}

\paragraph{Intuition.} \Cref{fig:logtan} draws \eqref{eq:logtan}. Plot the amplitude $\Mzero|v|$ against
$s=\ln|\tan(\alpha/2)|$: it is a $\operatorname{sech}$ bump of height $\Mzero\zeta$ centred at the Ernst angle, and a scan
samples it at $s=\ln|\tan(\Bone\alpha^{\rm nom}/2)|$. The MPME sequence is not spoiled, but by
\cref{thm:uv} each of its nine signals is $v$ times a function of $u$ alone, so all of them
report the same sample. The ratios between pathways give $u=\tanh(s-s_E)$, the offset of the
sample from the centre of the bump, and the overall scale then gives the height $\Mzero\zeta$.
Neither tells where the sample and the bump lie in absolute terms. Raising $\Bone$ moves the
sample to the right by $\ln[\tan(B_1^{+\prime}\alpha^{\rm nom}/2)/\tan(\Bone\alpha^{\rm nom}/2)]\approx\ln(B_1^{+\prime}/\Bone)$; shortening $\Tone$
by the right amount (about the square of the same factor, since $\zeta\approx\sqrt{\TR/2\Tone}$) moves the
bump by the same distance; and lowering $\Mzero$ by the factor restores the height. The whole
picture translates rigidly, and the scan cannot tell. In words: a larger flip angle in a tissue
with a larger Ernst angle (shorter $\Tone$) is indistinguishable from a smaller flip angle in a
tissue with a smaller Ernst angle.

A second scan samples the same bump at a different $s$. A $330^\circ$ pulse is $-30^\circ$ modulo
$360^\circ$, but its sample moves with $\Bone$ at the rate $\mathrm d\ln|\tan(\alpha_2/2)|/\mathrm d\ln\Bone=\alpha_2/\sin\alpha_2=-11.5$,
against $+1.01$ for the $15^\circ$ pulse. When the family translates the picture to keep scan 1,
the scan-2 sample therefore lands elsewhere on the bump (\cref{fig:logtan}a). A second small
angle such as $30^\circ$ moves at almost the same rate as the first ($1.05$), so its sample
translates with the bump and nearly preserves the data (\cref{fig:logtan}b). This is the
leading-order degeneracy of \cref{prop:scaling} for small-angle protocols.

For small angles, $\xi^2\approx\alpha^2\Tone/2\TR$ is the ratio of the saturation per pulse ($\propto\alpha^2$) to
the recovery per repetition ($\propto\TR/\Tone$), a useful mnemonic for the spoiled part of the signal.
The exact invariant uses half-angle tangents: along the family of \cref{tab:family}, the plain
ratio of the flip angle to the Ernst angle drifts from $1.0851$ to $1.0836$, while $\xi=1.0857$ is
constant.

\begin{figure}[!htbp]
\centering
\includegraphics[width=\linewidth]{ernst_logtan.png}
\caption{The single-scan degeneracy on the log half-angle axis (WM, $\TR=25$\,ms). Curves: the
amplitude $\Mzero|v|=\Mzero\zeta\operatorname{sech}(s-s_E)$ for three members of the family \eqref{eq:family}
($B_1^{+\prime}=0.8$, $1$, $1.25$); dotted lines: their Ernst angles $s_E$. Circles: the $15^\circ$ scan-1
sample of each member, at the same offset from its bump and the same height, hence identical
scan-1 data. Squares: the scan-2 sample. (a) For $330^\circ$ the scan-2 samples land at different
heights (and, for $B_1^{+\prime}=1.25$, with the opposite sign of $v$). (b) For $30^\circ$ they almost
coincide: a second small flip angle hardly separates the members.}
\label{fig:logtan}
\end{figure}

\paragraph{Relation to the spoiled gradient echo.} Let $m_z=(1-E_1)/(1-E_1\cos\alpha)$ be the longitudinal
magnetisation (per unit $\Mzero$) just before each pulse in the steady state of a spoiled gradient echo
with the same $\alpha$, $\TR$ and $\Tone$. Then
\begin{equation}
  v=\sin\alpha\;m_z,\qquad u=1-(1+\cos\alpha)\,m_z ,
  \label{eq:uvspoiled}
\end{equation}
since $1-u=(1-E_1)(1+\cos\alpha)/(1-E_1\cos\alpha)$. So $v$ is the spoiled (Ernst) signal: the transverse
magnetisation that each pulse creates from recovered longitudinal magnetisation (the
\emph{source}). $u\in[-1,1]$ is the factor by which one pulse-and-relaxation cycle carries existing
transverse magnetisation into the next repetition (the \emph{retention}, an effective cosine). In
the proof of \cref{thm:uv} they appear as the inhomogeneous term and the coefficient of the
transverse steady-state equation \eqref{eq:mperp_lin}. Unbalanced SSFP differs from the spoiled
sequence only in keeping the transverse magnetisation, and it sees $\alpha$ and $\Tone$ only through
this source and retention.

\paragraph{What one scan measures.} A scan yields 3 pathways $\times$ 3 echoes = 9 complex numbers
(18 real values). By \cref{thm:uv,prop:ernst} and \eqref{eq:signal}, they depend on the
parameters as listed in \cref{tab:onescan}.

\begin{table}[!htbp]
\centering
\caption{Information content of one scan.}
\label{tab:onescan}
\small
\begin{tabular}{ll}
\toprule
Feature of the data & Determined by\\
\midrule
amplitude ratios between pathways, $a_k/a_0=G_k(u)/G_0(u)$ & $u$, i.e.\ $\xi$ (and $E_2$)\\
overall scale, $\Mzero v=\Mzero\zeta\sin2\vartheta$ & $\Mzero\zeta$, given $\xi$\\
decay across the echoes of each pathway & $\Ttwo$ and $\Rtwop$ (hence $\Ttwos$)\\
phase change across echoes & $\dw$\\
common phase & $\phi_0$\\
\bottomrule
\end{tabular}
\end{table}

The 18 values therefore determine six quantities, $(\xi,\,\Mzero\zeta,\,\Ttwo,\,\Rtwop,\,\dw,\,\phi_0)$, and are
redundant with respect to them (which averages noise and allows model checking). The
deficiency is specific: the three unknowns $(\Mzero,\Bone,\Tone)$ enter only through the two
combinations $\xi$ and $\Mzero\zeta$ (equivalently $u$ and $\Mzero v$). More echoes or more pathways
from the same scan do not help, because every one of them depends on $(\alpha,\Tone)$ through the
same pair.

\paragraph{The data model for $(\Mzero,\Bone,\Tone)$.} Write $\xi_i$ and $A$ for the values of $\xi$
and $\Mzero\zeta$ that the data of scan $i$ determine. With $\alpha_i=\Bone\alpha_i^{\rm nom}$, the equations
that $(\Mzero,\Bone,\Tone)$ must satisfy are
\begin{equation}
  \frac{\tan(\Bone\alpha_i^{\rm nom}/2)}{\sqrt{\tanh(\TR/2\Tone)}}=\xi_i,\qquad
  \Mzero\sqrt{\tanh(\TR/2\Tone)}=A ,
  \label{eq:datamodel}
\end{equation}
equivalent to $u(\alpha_i,E_1)=U_i$ and $\Mzero v(\alpha_i,E_1)=V_i$ with $U_i=(\xi_i^2-1)/(\xi_i^2+1)$ and
$V_i=2A\xi_i/(1+\xi_i^2)$. Because all scans share $\TR$, the amplitude equation is the same for every
scan; only $\xi_i$ is new in each. For \emph{one} scan, \eqref{eq:datamodel} is two equations in three
unknowns, and its solutions form a one-parameter family. With the trial transmit scale
$B_1^{+\prime}$ as the parameter, $\tan(\alpha_E'/2)=\tan(B_1^{+\prime}\alpha^{\rm nom}/2)/\xi$, that is,
\begin{equation}
  \Tone'(B_1^{+\prime})=\frac{\TR}{2\,\operatorname{artanh}\bigl(\tan^2(B_1^{+\prime}\alpha^{\rm nom}/2)/\xi^2\bigr)},\qquad
  \Mzero'(B_1^{+\prime})=\Mzero\,\frac{\xi\,\zeta}{\tan(B_1^{+\prime}\alpha^{\rm nom}/2)} .
  \label{eq:family}
\end{equation}
In terms of $E_1$ this is $E_1'=(U+c')/(1+Uc')$ with $c'=\cos(B_1^{+\prime}\alpha^{\rm nom})$, the form used in
\cref{prop:singlescan}. A member exists whenever $|\tan(B_1^{+\prime}\alpha^{\rm nom}/2)|<|\xi|$, equivalently
$c'>-U$ (the Ernst angle must stay below $90^\circ$). Every member reproduces the scan exactly
(\cref{prop:singlescan}). With $x=B_1^{+\prime}\alpha^{\rm nom}$ and $\tau'=\TR/\Tone'$, differentiating \eqref{eq:family} gives
\begin{equation}
  \frac{\mathrm d\ln\Mzero'}{\mathrm d\ln B_1^{+\prime}}=-\frac{x}{\sin x},\qquad
  \frac{\mathrm d\ln\Tone'}{\mathrm d\ln B_1^{+\prime}}=-2\,\frac{x}{\sin x}\,\frac{\sinh\tau'}{\tau'},
  \label{eq:familyslopes}
\end{equation}
the exact forms of the factors $-1$ and $-2$ of \cref{prop:scaling}. For WM at $15^\circ$ they are
$-1.0115$ and $-2.023$, the slope measured in \cref{tab:singlescan}.

\paragraph{What a second flip angle adds.} Dividing the first equation of \eqref{eq:datamodel} for
$i=2$ by that for $i=1$ eliminates $\Tone$ and $\Mzero$:
\begin{equation}
  f(\Bone)\equiv\frac{\tan(\Bone\alpha_2^{\rm nom}/2)}{\tan(\Bone\alpha_1^{\rm nom}/2)}=\frac{\xi_2}{\xi_1}.
  \label{eq:secondscan}
\end{equation}
This is one equation in $\Bone$ alone, the half-angle analogue of the double-angle method. Its roots
are isolated, and $\Tone$ and $\Mzero$ then follow from \eqref{eq:family}. Its sensitivity to $\Bone$ is
the exact lever
\begin{equation}
  L=\frac{\mathrm d\ln|f|}{\mathrm d\ln\Bone}=\frac{\alpha_2}{\sin\alpha_2}-\frac{\alpha_1}{\sin\alpha_1}
  \qquad(\text{actual angles}),
  \label{eq:leverexact}
\end{equation}
$-12.53$ for $15^\circ/330^\circ$ and $+0.036$ for $15^\circ/30^\circ$ at $\Bone=1$ (\cref{sec:information}).

The data fix $\xi_i^2$ through $u_i$. Complex data with a phase $\phi_0$ shared by the two scans also fix
the relative sign of $v_1$ and $v_2$, hence the sign of $\xi_2/\xi_1$; magnitude data fix only
$|\xi_2/\xi_1|$. The aliases of \cref{prop:equiv} are the other roots of $f(B_1^{+\prime})=f(\Bone)$ (complex
data) or $|f(B_1^{+\prime})|=|f(\Bone)|$ (magnitude data). Their $\Bone$ values depend on the true $\Bone$ alone,
not on the tissue; the tissue enters only through the alias $\Tone'$, from \eqref{eq:family}, and
through admissibility ($\zeta'<1$). For the published protocol and true $\Bone=1$ ($f=-2.04$), the
magnitude alias is $B_1^{+\prime}=1.20$ ($f=+2.04$, scan 2 negated); it lies inside the physiological range and
is removed only by the shared phase (the FID-sign test of \cref{sec:magnitude}). The complex alias
is $B_1^{+\prime}=2.01$, which the bound $\Bone\le540/330$ excludes. Within the fitting range
$\Bone\in[\E^{-1},540/330]$, the complex equation has a single root for every true $\Bone$ between $6/11$ and
$1.592$, i.e.\ true $\alpha_2$ between $180^\circ$ and $525^\circ$. Because $f$ is the same function in every
voxel, a $\Bone$ shared by many voxels has the same aliases (\cref{sec:shared}).

\paragraph{A concrete example.} \Cref{tab:family} evaluates \eqref{eq:family} for one
white-matter voxel ($\Tone=850$\,ms, $\Mzero=1$, $\Bone=1$, scan 1 at $15^\circ$, $\TR=25$\,ms; $\xi=1.086$). As $B_1^{+\prime}$
runs from $0.8$ to $1.5$, the flip angle and the Ernst angle move together (their half-angle tangents
keep the ratio $\xi$). $\Tone'$ falls from $1334$ to $372$\,ms and $\Mzero'$ from $1.25$ to $0.66$. Every
member reproduces all nine complex scan-1 signals to machine precision, while the scan-2
($330^\circ$) signals differ by $41\%$ to $200\%$ (of the largest scan-2 signal) for every member except
the truth. The predicted ratio $\xi_2/\xi_1$ for a $330^\circ$ scan changes from $-10.6$ to $+12.1$ across the
family, whereas for a $30^\circ$ scan it changes only from $2.022$ to $2.082$. The second scan fixes
$\Bone$ when its angle moves differently with $\Bone$; more scan-1 data cannot.

\begin{table}[!htbp]
\centering
\caption{The single-scan solution family \eqref{eq:family} for one white-matter voxel (truth:
$\Bone=1$, $\Tone=850$\,ms, $\Mzero=1$). $B_1^{+\prime}\alpha$: actual scan-1 flip angle; $\alpha_E'$: Ernst angle of
$\Tone'$. Misfit: largest difference of the nine complex signals from the true ones, relative to
the largest true signal of that scan. Predicted $\xi_2/\xi_1$: the ratio $f(B_1^{+\prime})$ of
\eqref{eq:secondscan} that the member predicts for a second scan at $330^\circ$ or $30^\circ$; the data
have $f(1)$. Generated by \texttt{scripts/single\_scan\_family.py}.}
\label{tab:family}
\small\setlength{\tabcolsep}{4.5pt}
\begin{tabular}{ccccccccc}
\toprule
& & & & & \multicolumn{2}{c}{misfit} & \multicolumn{2}{c}{predicted $\xi_2/\xi_1$}\\
\cmidrule(lr){6-7}\cmidrule(lr){8-9}
$B_1^{+\prime}$ & $B_1^{+\prime}\alpha$ ($^\circ$) & $\alpha_E'$ ($^\circ$) & $\Tone'$ (ms) & $\Mzero'$ & scan 1 ($15^\circ$) & scan 2 ($330^\circ$) & $330^\circ$ & $30^\circ$\\
\midrule
0.80 & 12.0 & 11.06 & 1334 & 1.253 & $1\times10^{-15}$ & 0.69 & $-10.57$ & $2.022$\\
0.90 & 13.5 & 12.44 & 1052 & 1.112 & $2\times10^{-15}$ & 0.41 & $-5.18$ & $2.028$\\
1.00 & 15.0 & 13.83 & 850 & 1.000 & $2\times10^{-16}$ & 0.00 & $-2.04$ & $2.035$\\
1.10 & 16.5 & 15.21 & 701 & 0.908 & $3\times10^{-16}$ & 1.36 & $+0.18$ & $2.043$\\
1.20 & 18.0 & 16.60 & 587 & 0.831 & $5\times10^{-16}$ & 2.00 & $+2.05$ & $2.051$\\
1.50 & 22.5 & 20.76 & 372 & 0.662 & $1\times10^{-15}$ & 1.27 & $+12.14$ & $2.082$\\

\bottomrule
\end{tabular}
\end{table}

\subsection{Echo evolution and the voxel signal}

Pathway $k$ read at time $t$ after the pulse carries, for one isochromat, the off-resonance
phase $\dw(t+k\TR)$ (\cref{prop:offres} plus precession since the pulse). An intravoxel
frequency distribution with Lorentzian density $\ell(\delta)=(\Rtwop/\pi)/(\Rtwop{}^2+\delta^2)$
averages this to
\begin{equation}
  \int_{-\infty}^{\infty}\ell(\delta)\,\E^{i\delta\tau}\,d\delta=\E^{-\Rtwop|\tau|},
  \label{eq:lorentz}
\end{equation}
the characteristic function of the Cauchy distribution. Together with transverse relaxation
since the pulse, the noise-free signal of scan $i$, echo $j$, pathway $k$ is
\begin{equation}
  \boxed{\;S_{ijk}=\Mzero\,a_k\bigl(\Bone\alpha_i^{\rm nom},\Tone,\Ttwo,0\bigr)\,
  \E^{-t_{ijk}/\Ttwo}\,\E^{-\Rtwop|t_{ijk}+k\TR|}\,
  \E^{i(\phi_0+\dw(t_{ijk}+k\TR))}\;}
  \label{eq:signal}
\end{equation}
For $k\ge0$ the magnitude decays at $1/\Ttwo+\Rtwop=1/\Ttwos$; for $k<0$ and $t<|k|\TR$ at
$1/\Ttwo-\Rtwop$, a spin echo forming at $t=|k|\TR$. Echo times differ between pathways, and
$t+k\TR$ differs in sign between them, which is what separates $\Ttwo$ from $\Rtwop$.

\subsection{Closed forms for $k=0,-1$}

For the two lowest pathways the steady-state magnitudes have the closed forms of unbalanced
FISP and PSIF,
\begin{align}
  |a_0|&=\Bigl|\tan\tfrac{\alpha}{2}\Bigl[1-\frac{(E_1-\cos\alpha)(1-E_2^2)}{r}\Bigr]\Bigr|,\qquad
  |a_{-1}|=\Bigl|\frac{\tan\frac{\alpha}{2}}{E_2}\Bigl[1-\frac{(1-E_1\cos\alpha)(1-E_2^2)}{r}\Bigr]\Bigr|,
  \label{eq:fisp}\\
  r^2&=p^2-q^2,\quad p=1-E_1\cos\alpha-E_2^2(E_1-\cos\alpha),\quad q=E_2(1-E_1)(1+\cos\alpha).
  \nonumber
\end{align}
The PSIF formula in the literature is the echo amplitude at $t=\TR$, hence the factor $1/E_2$.
Dividing $p$ and $q$ by $1-E_1\cos\alpha$ shows that both brackets depend on $(\alpha,E_1)$ only
through $u$, and $\tan(\alpha/2)=v/(1-u)$, in agreement with \cref{thm:uv}.

\subsection{Imaging and noise}

Image formation $y_{ijkc}=\mathcal M_i\mathcal F(C_cS_{ijk})+n_{ijkc}$ (coil maps $C_c$, Fourier
transform $\mathcal F$, sampling $\mathcal M_i$) is linear. Voxel-wise, the coil-combined
complex images are modelled as $S_{ijk}+n_{ijk}$ with independent circular Gaussian noise of
standard deviation $\sigma$ per real and imaginary part; magnitudes are then Rician. The
published low-resolution scan 2 is modelled in \cref{sec:phantom} by truncation of $k$-space.

\subsection{Validation}

\begin{table}[!htbp]
\centering
\caption{Numerical checks of the forward model (all are unit tests).}
\label{tab:validation}
\small
\begin{tabular}{p{0.55\linewidth}l}
\toprule
Check & Agreement\\
\midrule
Isochromat solver vs truncated EPG ($K=400$), pathways $-3\ldots2$, four tissues incl.\ CSF & $<10^{-5}$ absolute\\
$|a_0|$, $|a_{-1}|$ vs closed forms \eqref{eq:fisp} & $10^{-6}$ relative\\
\Cref{thm:uv}: $a_k$ at $(\alpha',E_1')$ with equal $u$ equals $(v'/v)\,a_k$, six pathways, incl.\ $330^\circ$, $\dw\ne0$ & $10^{-9}$ relative\\
\Cref{prop:offres}: $a_k(\dw)=a_k(0)\E^{ik\dw\TR}$ & $10^{-10}$ absolute\\
Echo log-slopes $-(1/\Ttwo+\Rtwop)$ ($k\ge0$), $-(1/\Ttwo-\Rtwop)$ ($k=-1$) & exact\\
128 vs 16\,384 isochromats, CSF at $15^\circ$ and $330^\circ$ & $10^{-15}$ relative\\
\bottomrule
\end{tabular}
\end{table}
